\chapter{Failed, Abandoned and Underperforming Approaches}\label{ch:failures}

This chapter records every approach that failed, was abandoned, or did less than expected, together with
infrastructure failures that affected results or schedules. Each entry states:
\begin{itemize}
\item what was attempted and why;
\item what was found;
\item whether the explanation is \emph{verified} (checked by a further experiment or inspection) or a
\emph{hypothesis};
\item what changed as a result.
\end{itemize}
Entries are in chronological order.

\section{Method and modelling}

\subsection{F1. Reproducing the vanilla scores with the base checkpoint (2026-11-04)}\label{sec:fail-gate}
\begin{failbox}{}
\textbf{Attempt.} Llama-3.1-8B base, the checkpoint the README links, on the four rebuilt test sets.

\textbf{Why.} The paper names ``Llama-3.1-8B''. The vanilla gate ($\pm0.05$) must pass before fine-tuned gains can be
compared with the paper.

\textbf{Configurations tried} (\cref{tab:gen-vanilla-sql}). Six settings, tried on WikiSQL or more, against the
paper's 0.525:

{\small
\rowcolors{2}{rowshade}{white}
\begin{tabular}{@{}lll@{}}
\toprule\rowcolor{headshade}Setting & Model & WikiSQL result\\\midrule
Chat, no system message & Instruct & 0.552 (975 cases)\\
Chat, repo system message & Instruct & 0.519 (975 cases)\\
Raw completions, zero-shot & Base & 0.326\\
Raw completions, few-shot & Base & 0.685\\
Original-case tables, zero-shot & Base & 0.142\\
Original-case tables, few-shot & Base & 0.287\\
\bottomrule
\end{tabular}
}

\textbf{Result.} No base setting reproduced the paper on all four sets. Zero-shot met the gate only on BIRD; few-shot
met it on WikiTQ and BIRD.

\textbf{Against expectation.} We expected the README's checkpoint to reproduce its own paper.

\textbf{Explanation} (partly verified). The repository can query only through chat, which base Llama cannot serve.
Instruct through that chat call does reproduce the vanilla scores (\cref{sec:sql-vanilla}): this part is verified.
That the paper used Instruct is a hypothesis.

\textbf{Lesson and influence.} We stopped after six settings, because further combinations would eventually hit 0.525
by chance. The base gaps were recorded as unexplained. This later motivated the Instruct track.
\end{failbox}

\subsection{F2. The first validator (2026-11-04)}\label{sec:fail-validator}
\begin{failbox}{}
\textbf{Attempt.} A first 300-table, two-round TS run with our first comparator.

\textbf{Result.} It kept 42 and 31 pairs. A hand inspection of 50 kept and 24 rejected pairs, plus a re-run of all
160 \texttt{scala\_failed} pairs, found five defects:
\begin{enumerate}
\item SQL ran on all-text columns, so \texttt{MAX("Weight")} and \texttt{"n" < 50} compared strings;
\item a one-column NULL result was counted as a disagreement;
\item booleans did not match their 0/1 encoding;
\item Spark's \texttt{\textbackslash"} escaping was misread;
\item a fixed adapter name would have shared cached answers across arms.
\end{enumerate}

\textbf{Verified.} Each defect was reproduced on named cases (\texttt{gt\_05606}, \texttt{gt\_02180}, and others).
All 160 \texttt{scala\_failed} pairs exited with Spark code~2: 89 named a column that does not resolve and 51 called a
method the type lacks. None needed a validator change.

\textbf{Influence.} All five were fixed before the core study, and the first run was superseded
(\fpath{runs/old_test300_before_comparator_fix}). The corrected run's inspection of 50 kept and 50 rejected pairs
found that about 1 in 3 rejected pairs may hold usable SQL. That finding motivated Extensions~2 and~3.
\end{failbox}

\subsection{F3. The reduced-scope blueprint (abandoned 2026-11-04)}
\begin{failbox}{}
\textbf{Attempt.} Instruct with QLoRA, Pandas as the dual language, 2 rounds.

\textbf{Why abandoned.} The deadline was removed, so the scope-saving choices were reversed in favour of the paper's:
\begin{itemize}
\item LoRA: measured at 20.0~GB at 4,000 tokens, so it fits;
\item Spark-Scala: about 11~h per round estimated at the time, against 0.8~h for Pandas;
\item 3 rounds.
\end{itemize}
The model choice moved to base, by the README.

\textbf{Lesson.} Resource estimates in the blueprint were wrong in both directions. LoRA fit, but Spark was about 50
times slower per program than assumed. Measuring early let both reversals rest on data.
\end{failbox}

\subsection{F4. DPO on validator-labelled pairs (2026-11-07 to 2026-11-08)}
\begin{failbox}{}
\textbf{Attempt.} DPO with an SFT term (RPO), from the bestof4 adapter, on 397--438 (chosen, rejected) pairs per seed
(\cref{sec:method-ext2}).

\textbf{Why.} Rejected programs carry information about what not to write.

\textbf{Result.} $-0.0209$ against TS round~1, and below bestof4 on all 3 seeds (\cref{sec:ext2}). Five gentler variants
were tuned on dev under a rule fixed beforehand; none beat bestof4 there.

\textbf{Explanation} (partly verified).
\begin{itemize}
\item Verified: execution errors multiply (for example WikiTQ 15 $\to$ 58), and reward margins of 3.2--5.4 show strong
drift from the reference.
\item Verified on dev: the learning rate, not $\beta$, drives the damage.
\item Hypothesis: lowering the likelihood of plausible-but-wrong SQL also lowers nearby correct SQL.
\end{itemize}

\textbf{Influence.} DPO was dropped and reported as not helping. Effort moved from penalising rejected data to
recovering it.
\end{failbox}

\subsection{F5. Best-of-4 sampling: more data, little gain (2026-11-08)}
\begin{failbox}{}
\textbf{Result.} 52--62\% more validated questions, but only $+0.0064$ \Mviii{}, with seed differences of $+0.0282$
to $-0.0156$.

\textbf{Explanation} (hypothesis). The extra validated programs are mostly for questions that the greedy program
already answers, or nearly so.

\textbf{Influence.} Rounds 2--3 of bestof4 were put on hold. The same-data control was designed later to separate
volume from content.
\end{failbox}

\subsection{F6. Base Llama as its own repairer (2026-11-08)}\label{sec:fail-baserepair}
\begin{failbox}{}
\textbf{Attempt.} The first repair smoke test (12 pairs), then a pilot on 237 rejected pairs, with base as repairer
and the repository's correction phrase as prompt.

\textbf{Result.}
\begin{itemize}
\item In the smoke test, 6 of 10 greedy repairs copied the failed code exactly. Up to 3 retries at $T=0.7$ were added
in response.
\item In the pilot, base still copied the failed code in 85 of 169 Scala and 42 of 158 SQL repairs.
\item It recovered only 7 pairs (+17.1\%, against the required +30\%), so no-go.
\end{itemize}

\textbf{Explanation} (verified by the copy counts). A base model continues the text pattern rather than following a
repair request.

\textbf{Influence.} Instruct became the repairer (+56.1\%, go). This adds a second, stronger model to the base track,
which confounds repair with distillation (\cref{ch:discussion}).
\end{failbox}

\subsection{F7. Repair-TS on Instruct NL-to-SQL (2026-11-12 to 2026-11-14)}
\begin{failbox}{}
\textbf{Attempt.} The same repair-TS, with Instruct as both generator and repairer, 3 seeds $\times$ 3 rounds.

\textbf{Expectation.} Base had gained $+0.0421$, so we expected a positive effect.

\textbf{Result.} No effect at any round (round~3: $-0.0014$, $p=0.55$). Repaired pairs were worse than the same number
of extra first-pass pairs ($-0.0114$, $p<0.0001$).

\textbf{Explanation} (partly verified).
\begin{itemize}
\item Verified: recovery is lower than on base (+15--23\% against +52\% in round~1), and the largest rejected group,
\texttt{always\_empty}, is almost never recovered.
\item Hypothesis: when the generator and repairer are the same model, repair adds nothing the model does not already
know. On base, the stronger repairer added such knowledge.
\end{itemize}

\textbf{Influence.} Repair-TS was reframed as a conditional finding. The Scala scoring of the same runs then showed a
gain (\cref{sec:results-scala}).
\end{failbox}

\subsection{F8. Rounds on Instruct (2026-11-12)}
\begin{failbox}{}
\textbf{Expectation.} The paper says later rounds help, and base TS improved over rounds.

\textbf{Result.} Instruct TS is flat or declining: \ZSiv{} 0.6165, 0.6176, 0.6113. Round~3 against round~1 is
$-0.0083$ on \Mviii{} ($p=0.0012$). Seeds 1--2 peak at round~2.

\textbf{Explanation.} Not established. Instruct's round-1 adapter is already strong (\ZSiv{} 0.6165), and round sizes
stay small (486--529 pairs). Restarting from vanilla each round discards earlier rounds' data, so a later round sees
no more data than an earlier one. These are hypotheses.

\textbf{Influence.} Both the last-round and best-round means are reported against the paper.
\end{failbox}

\section{Pipeline and infrastructure}

\subsection{F9. Repository code that could not run as shipped}\label{sec:fail-repo}
\begin{failbox}{}
\begin{itemize}
\item \textbf{Evaluator.} \fpath{ns_evaluator.py} uses pandas APIs removed in pandas~3 (\texttt{applymap},
\texttt{errors="ignore"}) and placeholder paths. Fixed by compiling it with checked patches (\cref{lst:patches}).
\item \textbf{Sandbox.} The open-file limit of 256 stops Spark's JVM (\texttt{tzdb.dat (Too many open files)}; 0 of
32 programs ran). Fixed by raising it to 4,096 at load time, with no other limit changed.
\item \textbf{Training defaults.} The repository's defaults (\texttt{num\_few\_shot\_samples=10},
\texttt{sample\_method="random"}) raise an error in \texttt{generate\_train\_one\_dataset}. The authors must
therefore have trained either with \texttt{random-Gen} or with no few-shot examples; this is not recorded. We train
zero-shot prompts only, which is an unavoidable difference.
\end{itemize}
\end{failbox}

\subsection{F10. U+2028 in training data stopped a run (2026-11-05)}
\begin{failbox}{}
\textbf{What happened.} The no-validation arm stopped when round-2 SQL training started. Our JSONL reader used
\texttt{str.splitlines}, which also splits on U+2028, and 2 of 2,863 rows contained that character.

\textbf{Verified fix.} Split only on \texttt{\textbackslash n}, applied to all three readers, with a regression test.
The arm restarted from its saved round-2 data. Table-Specialist was unaffected (its files hold no U+2028).
\end{failbox}

\subsection{F11. Parallel BIRD scoring undercounted 3 cases (found 2026-11-11, fixed 2026-11-12)}\label{sec:fail-bird}
\begin{failbox}{}
\textbf{Symptom.} Greedy BIRD scores in the self-consistency runs were exactly 3 cases (0.0082) below the baseline
scorer, on every BIRD cell.

\textbf{Verified cause.} The 3 cases return 20,198--49,729 rows and hit the 5~s SQL limit while 16 scoring threads
run. The baseline scorer uses one thread.

\textbf{Fix.} BIRD is scored on one thread. The 14 base BIRD cells were re-scored under new identifiers
(\texttt{final\_sct10seq\_*}), keeping the old files. Greedy BIRD now matches the baseline exactly.
\end{failbox}

\subsection{F12. A void NL-to-Scala smoke test (2026-11-14)}
\begin{failbox}{}
\textbf{What happened.} The first smoke test's \texttt{tmux} command did not source the sandbox settings, so all 132
programs failed in 0.02~s with ``Set TABLE\_SPECIALIST\_SPARK\_IMAGE'', and those failures were cached.

\textbf{Fix.}
\begin{itemize}
\item The scorer calls \texttt{ensure\_sandbox} before running anything, so a missing setting stops the run instead of
being cached.
\item The cache moved to a new directory; the old one, holding only those 132 records, is never read.
\item The rerun used a new tag, so the void results were not overwritten.
\end{itemize}
\end{failbox}

\subsection{F13. Error-detection smoke test: a collapsed validator (2026-11-15)}
\begin{failbox}{}
\textbf{What happened.} Round~2 of the smoke run first failed because its vLLM server was not ready within 20 minutes
during the machine stall. The rerun completed, but its 5-step adapter answered NULL for every column. As validator it
rejected every injected error, and as a model it predicted nothing (F1 0.0).

\textbf{Reading.} A smoke-scale artefact of 5 training steps on 18 examples, flagged as the risk to watch: a validator
that never flags errors would leave later rounds only clean columns to train on.

\textbf{Verified at full scale.} The round-1 adapters still flag errors: seed~1 predicts on 187 of 1,126 ST columns.
\end{failbox}

\subsection{F14. Machine stall and stalled timeouts (2026-11-14, 19:37--03:03)}
\begin{failbox}{}
\textbf{What happened.} Nothing progressed for about 7.5~h. The journal shows systemd-journald restarting on an
uncleanly shut-down journal, and uptime shows no reboot. The likely cause, not certain, is that the Windows host paused
the WSL virtual machine.

\textbf{Effect.} 36 Spark programs in flight were cached as ``timed out'' after 26,891--26,955~s, although the limit is
120~s. These were the only timeouts among 1,758 cache records at the time.

\textbf{Fix.} A cached timeout longer than 600~s is treated as missing and re-run, with the result stored beside the
original record. The scorer holding those records in memory was stopped before writing any cell, and relaunched.
\end{failbox}

\subsection{F15. A quadratic comparison (found and fixed 2026-11-15)}\label{sec:fail-quadratic}
\begin{failbox}{}
\textbf{Symptom.} Writing one TS BIRD Scala cell took about 43~minutes, against about 1 minute for vanilla.

\textbf{Verified cause.} The order-insensitive comparison rebuilt two row-count differences for every row: quadratic in
the result size, about 10~minutes for the 49,729-row BIRD gold. TS triggers it often, because its results have the
gold's row count.

\textbf{Fix.} The value is computed once. All 51 cells written so far (45,930 cases) were rescored read-only with the
new code: 0 differences in any verdict, in 149~s in total. A new test checks a 20,000-row comparison runs in under 5~s.
Published NL-to-SQL corrected scores are unaffected for the same reason.
\end{failbox}

\subsection{F16. Instruct smoke tests keep very few pairs (2026-11-10)}
\begin{failbox}{}
\textbf{What happened.} The first Instruct TS smoke test (20 tables) kept 0 of 20 pairs (always-empty 11, disagree 4,
Scala failed 4, SQL failed 1), so training had no examples and the run exited with an error. A 60-table rerun kept 4
and 3 pairs (5--7\%).

\textbf{Reading.} Possibly chance at this size (base kept 14.7\% in round~1). Always-empty at 55\% was flagged to
watch.

\textbf{Verified at full scale.} Instruct keeps 10.7--11.1\% in round~1, and \texttt{always\_empty} is indeed its
largest rejection reason (\cref{tab:gen-validation}).
\end{failbox}

\subsection{F17. R repair feedback lost the cause of the error (found and fixed 2026-11-16)}\label{sec:fail-r-feedback}
\begin{failbox}{}
\textbf{Symptom.} In the NL-to-R smoke test, repair recovered 0 of 26 rejected pairs. For 4 of the 6 R programs that
failed, the feedback shown to the model was only \texttt{Error in filter():} or \texttt{Error in summarise():}.

\textbf{Verified cause.} \texttt{dplyr} (through \texttt{rlang}) prints the cause on the lines after the first, for
example \texttt{Caused by error: ! object 'Access' not found}. Our extractor kept only indented continuation lines,
which is R's base-error format, so the cause was dropped.

\textbf{Fix.} A message is now the error line plus every following line up to the backtrace, capped at 300
characters. A unit test uses a real dplyr error. Recomputed on the smoke test's six failures, every message carries its
cause (e.g.\ \texttt{could not find function "year"}). Only the repair arm reads this text, so TS-R was unaffected; the
full repair-TS-R runs started with the fix and recovered 84--126 pairs per round.
\end{failbox}

\subsection{F18. One experiment took another's GPU (2026-11-16)}
\begin{failbox}{}
\textbf{Symptom.} TS-R seed~1 round~2 took about 1.8~h longer than round~1, while GPU~1 sat idle.

\textbf{Verified cause.} The GPU lock takes the first free GPU, trying GPU~0 first. TS-R released GPU~0 for its
11-minute, CPU-only validation step; the running repair-ED started its next step in that window, took GPU~0 and held it
for about 2~h.

\textbf{Fix.} NL-to-R can keep its GPU through validation (\texttt{TS\_R\_HOLD\_GPU=1}), and repair-ED can be limited to
given GPUs (\texttt{TS\_ED\_GPUS}). TS-R was restarted with the setting during a cached generation step, losing nothing.
No result was affected: each step is self-contained; only time was lost.
\end{failbox}
